\documentclass[11pt]{article}
\usepackage{fontspec}
\usepackage[french]{babel}
\usepackage{amsmath}
\usepackage{amssymb}
\usepackage{amsthm}
\usepackage{mathrsfs}
\usepackage{tikz-cd}
\usepackage{enumitem}

\frenchsetup{AutoSpacePunctuation=false}

\newcommand{\printedpage}[1]{\marginpar{\footnotesize p.~#1}}
\newcommand{\msp}[1]{\marginpar{\footnotesize [#1]}}

\newcommand{\HOT}{\mathrm{HOT}}
\newcommand{\Hot}{\mathrm{Hot}}
\newcommand{\Cat}{\mathrm{Cat}}
\newcommand{\Ob}{\mathrm{Ob}}
\newcommand{\hot}{\mathrm{hot}}
\newcommand{\uHom}{\underline{\mathrm{Hom}}}
\newcommand{\uFib}{\underline{\mathrm{Fib}}}
\newcommand{\uFibsc}{\underline{\mathrm{Fibsc}}}
\newcommand{\uCh}{\underline{\mathrm{Ch}}}
\newcommand{\circled}[1]{\tikz[baseline=(c.base)]\node[draw,rounded corners=4pt,inner sep=1.5pt](c){#1};}

\begin{document}

\printedpage{79}
% head: A. Grothendieck : Dérivateurs, Ch. VI, Édition des Universités de Montpellier II et Paris VII
% foot: Version 9 novembre 2001

donc pour $S$, $S'$ données dans Cat, définir
\[ \uHom(S',S)^o \longrightarrow \uHom(\HOT_W(S),\HOT_W(S')). \]
Cela revient surtout, pour
\[ f,g : S' \rightrightarrows S \]
deux flèches de $S'$ dans $S$, et un homomorphisme de foncteurs
\[ u : f \longrightarrow g, \]
à décrire un homomorphisme de foncteurs
\[ u^* : g^* \longrightarrow f^* \]
entre $\HOT_W(S)$ et $\HOT_W(S')$. C'est d'ailleurs chose faite, dans le cadre général des $(\mathcal{M},W)$, quand on décrit $\HOT_W(S)$ et $\HOT_W(S')$ par localisation de $\mathrm{Hom}(S^o,\Cat)$ et $\mathrm{Hom}(S'^o,\Cat)$. Ces catégories s'interprètent respectivement comme $\uFibsc\,S$ et $\uFibsc\,S'$, donc comme des catégories de catégories sur $S$ resp. $S'$, on voudrait donc interpréter $u^*$ en termes de telles catégories, et tant qu'à faire, pour des catégories $X$ sur $S$ aussi générales que possible. Nous nous
\msp{page 129}
attendrons alors à devoir faire des hypothèses, soit sur $X$, soit sur $f$ et sur $g$, qui assurent que les images de $\hot_{W,S}(X)$ par $f^*$ et $g^*$ ne sont autres que $\hot_{W,S}(f^*X)$, $\hot_{W,S}(g^*X)$ respectivement, où on pose
\[ f^*(X) = X \times_S (S',f), \qquad g^*(X) = X \times_S (S',g). \]
Donc on devra sans doute se placer dans l'une des hypothèses suivantes:
\begin{enumerate}
\item[a)] $X$ est lisse sur $S$.
\item[b)] $X$ est Hot-fibrée sur $S$.
\item[c)] $f$ et $g$ sont propres (et sous réserve de la validité de la conjecture page 122).
\end{enumerate}
On s'attend alors à ce que $u^*$ soit associé à un homomorphisme
\[ u^* : g^*(X) \longrightarrow f^*(X) \qquad \textit{dans } \Cat/S' \]
défini en termes de $u$. Y a-t-il moyen de le définir bel et bien, sans supposer $X$ au moins fibrée sur $S$? Prenons $S'=e$, donc $f$, $g$ correspondant à des objets $s,t\in\Ob\,S$, et $u$ à une fléche % [sic] fléche -> flèche, wrong accent
\[ u : s \longrightarrow t, \]
les images inverses de $X/S$ par $f$, $g$ sont les fibres $X_s$, $X_t$, et on cherche
\[ u^* : X_t \longrightarrow X_s. \]
On n'a aucune chance d'en trouver canoniquement, si on ne suppose $X$ fibrée sur $S$! Donc ce ne sont pas a), b) ni c) qui peuvent nous servir, mais unique%
\printedpage{80}%
% head: A. Grothendieck : Dérivateurs, Ch. VI, Édition des Universités de Montpellier II et Paris VII
% foot: Version 9 novembre 2001
\msp{page 130}%
ment l'hypothèse $X\in\Ob\,\uFib\,S$. Donc on ne peut espérer trouver une description de $u^*$ entre $\HOT_W(S)$ et $\HOT_W(S')$ en termes d'opérations sur les catégories sur $S$, $S'$, comme suggéré à l'instant, qu'en se bornant à des `modèles' $X$ qui sont dans $\uFib\,S$, tout au moins. Et je ne doute pas en effet que ça marche sur $\uFib\,S$ pas moins bien que directement sur $\uFibsc\,S$, en termes de $\uHom(S^o,\Cat)$.

Pour les axiomes essentiels sur les dérivateurs, je renvoie à I p. 67-68 et p. 111-112. On a vu que Der 1, \circled{Der 2} (\footnote[66]{Canulé pour \underline{Der 2}, je ne sais si c'est vrai même en supposant $W$ \textit{fortement} saturé, i.e. tel que $\hot_W(f)$ isomorphisme $\Longrightarrow f\in W$.}) sont satisfaits pour tout prédérivateur défini par un couple $(\mathcal{M},W\subseteq\mathrm{Fl}\,\mathcal{M})$, donc c'est OK ici pour $\HOT_W$. Il s'ensuit que c'est OK aussi pour $\HOT_W^{\mathrm{lc}}$, vu la caractérisation des $\HOT_W^{\mathrm{lc}}(S)$ en tant que sous-catégories de $\HOT_W(S)$. Pour Der 1, ça marche-t-il même pour des sommes quelconque de catégories, i.e. a-t-on % [sic] quelconque -> quelconques, agreement with sommes; Canulé (footnote 66) -> Annulé ?, not a French word
\[ \HOT_W(\coprod_i S_i) \xrightarrow{\sim} \prod_i \HOT_W(S_i)\,? \]
On a en tous cas
\[ \uFib(\overbrace{\coprod\limits_i S_i}^{S}) \xrightarrow{\sim} \prod_i \uFib(S_i), \]
et le localiseur $W_S^{\mathrm{f}}$ dans $\uFib(S)$ s'identifie
\msp{page 131}
bel et bien au produit des localiseurs $W_{S_i}^{\mathrm{f}}$ dans les $\uFib\,S_i$. Mais dans le cas d'une famille peut-être infinie de catégories $\mathcal{F}_i$, munies de localiseurs $W_i$, si on pose
\[ \mathcal{F} = \prod\mathcal{F}_i, \qquad W = \prod W_i \subseteq \mathrm{Fl}(\mathcal{F}), \]
le foncteur canonique
\[ \mathcal{F}W^{-1} \longrightarrow \prod_i(\mathcal{F}_iW_i^{-1}) \]
est-il toujours une équivalence? Même dans \texttt{[`Même dans' barré]} le cas $I$ fini il faut (il me semble) un argument assez soigneux, soit par la propriété universelle par récurrence sur card\,$I$ (on est ramené au cas de deux foncteurs), soit à coups de catégories de chemins, en notant que
\[ \uCh_\infty^W(\mathcal{F};(s_i),(t_i)) \xrightarrow{\sim} \prod_i \uCh_\infty^{W_i}(\mathcal{F}_i;s_i,t_i), \tag*{$(*)$} \]
et en passant aux $\pi_0$, qui donne
\[ \mathrm{Hom}_W((s_i),(t_i)) \xrightarrow{\sim} \prod_i \mathrm{Hom}_{W_i}(s_i,t_i). \]
Mais la formule $(*)$, visiblement, ne marche que pour $I$ fini. Dans le cas général, on trouve la sous-catégorie $\prod_i^!$ de $\prod_i$ formée des systèmes $(c_i)_{i\in I}$ de chemins qui sont de longueur bornée.

Donc la question de la validité de Der 1 sans hypothèse de finitude sur l'ensemble d'indices $I$
\printedpage{81}
% head: A. Grothendieck : Dérivateurs, Ch. VI, Édition des Universités de Montpellier II et Paris VII
% foot: Version 9 novembre 2001
\msp{page 132}
reste en suspens pour l'instant. (Je soupçonne que c'est vrai, mais il faudrait, pour le prouver, tenir une description des flèches dans $\HOT_W(S)$ par des chemins de longueur bornée, disons de longueur deux ou trois. C'est une question intéressante, mais que je réserve pour plus tard, si le besoin m'y contraint $\ldots\,$)

Plus importante est la question de la validité de Der 3 et de ses variantes (I, p. 111-112). Le minimum qu'on demande d'un dérivateur $\mathbf{D}$, c'est que si $S'\longrightarrow S$ dans Diag est une Hot-équivalence, ce soit une $\mathbf{D}$-équivalence, i.e. (dans le cas actuel) que le foncteur
\[ f^* : \HOT_W(S) \longrightarrow \HOT_W(S') \]
soit pleinement fidèle sur la sous-catégorie des objets \textit{constants} dans $\HOT_W(S)$. Cette condition devient plus tangible quand on dispose déjà soit des foncteurs cohomologie $H^\bullet_{\mathbf{D}}(S,-)$, $H^\bullet_{\mathbf{D}}(S',-)$, soit de l'homologie $H_\bullet^{\mathbf{D}}(S,-)$, $H_\bullet^{\mathbf{D}}(S',-)$. En termes de cette dernière, elle signifie que
\[ p_!p^*(a) \xrightarrow{\sim} p'_!p'^*(a) \qquad \text{\texttt{[est un]} isomorphisme,} \]
\msp{page 133}
pour tout $a$ dans $\mathcal{A}=\mathbf{D}(e)$. (Où $p:S\longrightarrow e$, $p':S'\longrightarrow e$ sont les morphismes structuraux sur la catégorie ponctuelle.) On a vu la formule heuristique générale
\[ p_!\underbrace{p^*(a)}_{a_S} = \hot(S)\times a, \]
où $\hot\,S$ est la classe de $S\in\Ob\,\Cat$ dans $\Hot=(\Cat)W_\infty^{-1}$, et où le produit $\times$ est une opération externe de Hot sur $\mathcal{A}=\mathbf{D}(e)$
\[ \Hot\times\mathcal{A} \longrightarrow \mathcal{A}\ . \]
Dans le cas actuel, $\mathcal{A}=\Hot$, et l'opération
\[ \Hot\times\Hot \longrightarrow \Hot \]
ne doit être autre que le foncteur produit cartésien, lequel est déduit du foncteur similaire dans Cat. Quand on dispose de ce formalisme, le résultat cherché pour un $f:S'\longrightarrow S$ se réduit simplement à dire que si $f$ est dans $W$, alors il en est de même de $f\times\mathrm{id}_X$ pour tout $X$ dans Cat (en prenant $X$ tel que $\hot(X)\simeq a$). Or là c'est un fait bien connu, résultant des axiomes mis sur $W$.

J'ai donc l'impression que je n'ai aucune chance de prouver Der 3a$'$ (p. 111), ni même Der 3a) (cas où $\mathrm{S}=\mathrm{e}$), sans avoir une maîtrise du formalisme des $p_!$. Et j'ai
\printedpage{82}
% head: A. Grothendieck : Dérivateurs, Ch. VI, Édition des Universités de Montpellier II et Paris VII
% foot: Version 9 novembre 2001
\msp{page 134}
bon espoir que cetter maîtrise va s'obtenir à bon compte, alors que la question des $p_*$ semble beaucoup plus délicate. % [sic] cetter -> cette, misspelling

Je vais donc me tourner d'abord vers l'axiome Der 4$'$ d'existence de
\[ f_! : \HOT_W(S') \longrightarrow \HOT_W(S), \]
comme adjoint à gauche de $f^*$. Je vais essayer de construire ad-hoc un foncteur $f_!$, puis les deux flèches d'adjonction
\[ \begin{aligned} \mathrm{id} &\quad\longrightarrow\quad f^*f_! \\ \mathrm{id} &\quad\longleftarrow\quad f_!f^*, \end{aligned} \]
et enfin, de vérifier les deux compatibilités pertinentes.

L'idée géométrique `évidente', c'est pour $X'\in\Cat/S'$, représentant $\xi'\in\Ob\,\HOT_W(S')$, de définir $f_!(\xi')$ par
\[ f_!(\hot_{S',W}(X')) = \hot_{S,W}(X'), \]
où dans le second membre, $X'$ est regardé comme objet de $\Cat/S$, grâce à $f:S'\longrightarrow S$ et au foncteur
\[ f_! : \Cat/S' \longrightarrow \Cat/S \]
qu'il induit. Le seul problème, c'est de voir si ce foncteur est compatible avec les
\msp{page 135}
localiseurs $W_{S'}^{\mathrm{d}}$, $W_S^{\mathrm{d}}$, et sinon, s'il l'est tout au moins sur la sous-catégorie $\uFib\,S'$ ou $\uFibsc\,S'$, qui suffit pour calculer $\HOT_W(S')$ par localisation. Il faut donc voir si un $X'_1\longrightarrow X'_2$ dans $\Cat/S'$, qui est dans $W_{S'}^{\mathrm{d}}$, i.e. qui est dans $W$ et le reste par tout changement de base de la forme $s'\backslash S'\longrightarrow S'$, plus généralement pour tout co-isomorphisme local, est aussi dans $W_S^{\mathrm{d}}$, i.e. si les changements de base de la forme $s\backslash S\longrightarrow S$ le transforment en $W$-équivalences (grossières). Or les co-isomorphismes locaux sont stables par \textit{tout changement de base} (cf. I, p. 55 dans le cas dual des isomorphismes locaux), donc on gagne!
\[
\begin{tikzcd}
\overline{S'} = s\backslash S' \arrow[r] \arrow[d] & S' \arrow[d] \\
\overline{S} = s\backslash S \arrow[r] & S
\end{tikzcd}
\quad .
\]
Je vais quand même le vérifier encore. Si
\[
\begin{tikzcd}
\overline{S'} \arrow[r] \arrow[d] & S' \arrow[d] \\
\overline{S} \arrow[r] & S
\end{tikzcd}
\]
\printedpage{83}
% head: A. Grothendieck : Dérivateurs, Ch. VI, Édition des Universités de Montpellier II et Paris VII
% foot: Version 9 novembre 2001
est un carré cartésien dans Cat, on a dit dans loc. cit., essentiellement, que pour $\overline{s'}$ dans $\overline{S'}$, ayant des images $s'$, $\overline{s}$, $s$ dans $S'$, $\overline{S}$, $S$ respectivement, le carré des colocalisés
\[
\begin{tikzcd}
\overline{s'}\backslash\overline{S'} \arrow[r] \arrow[d] & s'\backslash S' \arrow[d] \\
\overline{s}\backslash\overline{S} \arrow[r] & s\backslash S
\end{tikzcd}
\]
est encore
\msp{page 136}
cartésien. Je viens encore le de revérifier à l'instant, pour les objets tout au moins. Dire que $\overline{S}\longrightarrow S$ est un isomorphisme local, signifie que les $\overline{s}\backslash\overline{S}\longrightarrow s\backslash S$ sont \texttt{[des]} isomorphismes. Mais le carré cartésien précédent montre qu'alors les $\overline{s'}\backslash\overline{S'}\longrightarrow s'\backslash S'$ sont aussi \texttt{[des]} isomorphismes, on gagne! % [sic] le de revérifier -> de le revérifier, word order

\textbf{Théorème 1.} \textit{Pour toute flèche $f:S'\longrightarrow S$ dans} Cat\textit{, le foncteur induit}
\[ f_! : \Cat/S' \longrightarrow \Cat/S \]
\textit{est compatible avec les localiseurs $(W_{S'}^{\mathrm{d}},W_S^{\mathrm{d}})$, donc dualement avec les localiseurs $(W_{S'}^{\mathrm{g}},W_S^{\mathrm{g}})$, et trivialement avec les localiseurs $(W_{S'},W_S)$. Il induit donc un diagramme commutatif de foncteurs sur les localisés}
\[
\begin{tikzcd}
\HOT_W(S') \arrow[r, "\rho_{S'}"] \arrow[d, "f_!"] & \HOT_W^{\mathrm{lc}}(S') \arrow[d, "f_!^{\mathrm{lc}}"] & \HOT_W^o(S') \arrow[l, "\sigma_{S'}"'] \arrow[d, "f_!^o"'] \\
\HOT_W(S) \arrow[r, "\rho_S"] & \HOT_W^{\mathrm{lc}}(S) & \HOT_W^o(S). \arrow[l, "\sigma_S"']
\end{tikzcd}
\]

\textbf{Définition 1.} \textit{On a défini ainsi les foncteurs verticaux $f_!$, $f_!^{\mathrm{lc}}$, $f_!^o=(f^o)_!$, appelés foncteurs image directe homologique, ou foncteurs d'intégration pour $f$ (et les types d'homotopie relatifs).}

\msp{page 137}
On a obtenu en même temps, gratuitement, la compatibilité de ces foncteurs aux foncteurs \texttt{[de]} localisation $\rho_{S'}$, $\sigma_{S'}$ pour $\HOT(S')$, et $\rho_S$, $\sigma_S$ pour $\HOT(S)$. Par contre, on ne s'attendra pas qu'ils commutent avec les foncteurs pleinement fidèles d'inclusion des $\HOT_W^{\mathrm{lc}}$ dans $\HOT_W$. En effet, si on ne suppose pas $f$ dans \underline{Fibhot} (\footnote[67]{Ce cas doit être regardé avec soin ultérieurement.}), il est tout à fait immoral de s'attendre à ce que $f_!$ dans $\HOT_W(S')$ transforme types localement constants en itou. La situation est donc exactement inverse de celle pour le foncteur $f^*$, qui était compatible avec les foncteurs d'inclusion $i_S$, $i_{S'}$ et $j_S$, $j_{S'}$ - c'était à tel point trivial qu'il n'a pas été question même de le dire, voire en faire une proposition! La situation est

\end{document}
